\section*{Nomenclatura}

\begin{tabbing}
$A$\qquad\qquad \= Área da seção transversal (m\textsuperscript{2})\\
$b_w$ \> Largura da seção transversal do tabuleiro (m)\\
$CIV$ \> Coeficiente de impacto vertical\\
$d$ \> Diâmetro da longarina circular (m)\\
$esp$ \> Espaçamento entre longarinas (m)\\
$E_{0,\mathrm{ef}}$ \> Módulo de elasticidade efetivo (kPa)\\
$f_{m,d}$ \> Resistência de cálculo à flexão (kPa)\\
$f_{v0,d}$ \> Resistência de cálculo ao cisalhamento paralelo às fibras (kPa)\\
$f_m(\mathbf{x})$ \> $m$-ésima função objetivo\\
$g_j(\mathbf{x})$ \> $j$-ésima restrição de desigualdade\\
$h$ \> Altura da seção transversal do tabuleiro (m)\\
$I$ \> Momento de inércia da seção transversal (m\textsuperscript{4})\\
$k_{mod}$ \> Coeficiente de modificação\\
$L$ \> Vão teórico da longarina (m)\\
$M_{gk}$ \> Momento fletor característico devido à ação permanente (kN$\cdot$m)\\
$M_{qk}$ \> Momento fletor característico devido à ação variável (kN$\cdot$m)\\
$M_{sd}$ \> Momento fletor solicitante de cálculo (kN$\cdot$m)\\
$V_{sd}$ \> Esforço cortante solicitante de cálculo (kN)\\
$W$ \> Módulo resistente da seção transversal (m\textsuperscript{3})\\
$\mathbf{x}$ \> Vetor de variáveis de projeto\\
$\gamma_g,\ \gamma_q$ \> Coeficientes parciais das ações permanentes e variáveis\\
$\gamma_{wf},\ \gamma_{wc}$ \> Coeficientes de minoração da resistência (flexão e cisalhamento)\\
$\delta_{\mathrm{tot}}$ \> Deslocamento total no meio do vão (m)\\
$\delta_{\mathrm{lim}}$ \> Deslocamento limite admissível (m)\\
$\varphi$ \> Coeficiente de fluência\\
$\psi_2$ \> Fator de combinação quase permanente\\
$\sigma_{Sd},\ \sigma_{Rd}$ \> Tensão normal solicitante e resistente de cálculo (kPa)\\
$\tau_{Sd},\ \tau_{Rd}$ \> Tensão de cisalhamento solicitante e resistente de cálculo (kPa)\\
$\rho$ \> Fator de perturbação da variável de projeto (robustez)
\end{tabbing}
